\documentclass[10pt]{article}
\usepackage[a4paper,margin=20mm]{geometry}
\usepackage{amsmath,amssymb,booktabs,array,microtype}
\usepackage[hidelinks]{hyperref}
\usepackage{xcolor}

\hypersetup{pdftitle={Fast Complete BN254 Scalar Multiplication},pdfauthor={Kriterion research submission}}

\definecolor{ink}{HTML}{17212B}
\definecolor{accent}{HTML}{145A7A}
\definecolor{light}{HTML}{EAF3F7}
\color{ink}
\setlength{\parindent}{0pt}
\setlength{\parskip}{3pt}
\renewcommand{\arraystretch}{1.15}

\title{\vspace{-16mm}\textbf{Fast Complete BN254 Scalar Multiplication}\\
\large Twelve optimisations for ArgoMAC, with specification, security, and speedup}
\author{Kriterion research submission}
\date{6 September 2026}

\begin{document}
\maketitle

\begin{abstract}
This report presents twelve changes to the ArgoMAC garbled circuit for BN254 G1 scalar multiplication.
Each change has its own section.
Each section states whether the change alters the specification, gives the construction detail, gives the security analysis, and reports the measured speedup.
Three changes are in the Lean construction: one shared garbling pass, a proved 91-digit recoding, and complete Renes--Costello--Batina addition.
Nine changes are in the Rust implementation.
The Lean benchmark median falls from 989 ms to 231 ms.
The Rust protocol encode mean falls from 81.309 ms to 68.190 ms.
Lean 4.24 checks correctness, the digit bound, the complete law, the topology, and the Lamport label order without a new axiom.
The adaptive-privacy theorem is not complete; Section~\ref{sec:formal} states its exact status.
\end{abstract}

\section{Summary}

Table~\ref{tab:summary} lists every change.
The Lean column refers to the challenge construction and its proof.
The Rust column refers to the BaBe implementation.
A specification change is a change to the paper construction, the wire format, the randomness derivation, or the evaluator acceptance rule.
A change without a specification change computes the same function with the same public values.

\begin{table}[h]
\centering
\small
\begin{tabular}{llllr}
\toprule
\S & Change & Layer & Spec change & Measured effect \\
\midrule
\ref{sec:o1} & Shared garbling pass & Lean & No & Lean 989 to 160 ms \\
\ref{sec:o2} & Four-corner 91-digit recoding & Lean, Rust & Yes: digit count & One fewer output row \\
\ref{sec:o3} & Complete RCB addition & Lean, Rust & Yes: row law, decode & Correctness repair; Lean 160 to 231 ms \\
\ref{sec:o4} & Parallel rows with fixed RNG streams & Rust & Yes: randomness derivation & Row 15.934 to 3.845 ms \\
\ref{sec:o5} & Fixed AES window view reuse & Rust & No & Within noise \\
\ref{sec:o6} & SHA-256 prefix reuse & Rust & No & Row 3.841 to 3.754 ms \\
\ref{sec:o7} & Specialized hash-to-field reduction & Rust & No & One worker 20.858 to 19.906 ms \\
\ref{sec:o8} & Four-limb field XOR & Rust & No & Row 3.713 to 3.390 ms \\
\ref{sec:o9} & Direct hash-block reduction & Rust & No & Row 3.507 to 3.351 ms \\
\ref{sec:o10} & One-block KDF input & Rust & Yes: KDF layout, version 5 & Row 3.375 to 3.081 ms \\
\ref{sec:o11} & Paired KDF digests & Rust & Yes: KDF layout, version 6 & Row 3.121 to 2.876 ms \\
\ref{sec:o12} & Noncanonical coordinate rejection & Rust & Yes: acceptance rule & None; hardening \\
\bottomrule
\end{tabular}
\caption{All changes. Each Rust row time is one adjacent A/B pair from that change's own benchmark run, so the baselines drift between rows.}
\label{tab:summary}
\end{table}

All Rust times use an Apple M5 with ARMv8 AES, Rust 1.94.0, and 100 Criterion samples.
Row times use ten Rayon workers unless the text says one worker.
All Lean times are 21-run medians of the direct-table benchmark on the same machine.

\section{Shared garbling pass}
\label{sec:o1}

\subsection{Specification change}
None.
The logical circuit, the public table, the fixed-key oracle indices, the topology, the input labels, and the assumptions do not change.

\subsection{Detail}
The baseline executes the same pure garbling operations six times in its dominant layer.
It computes each bit-adaptor pair twice.
It computes each point row three times.
The shared pass computes each bit pair once and each point row once.
It then reuses the values.
The fixed-permutation application count falls from 6,322,060 to 1,057,910 for the original 92-row, nine-adaptor circuit.

\subsection{Security analysis}
Lean proves that both changed functions equal their baseline definitions for all inputs.
Thus, the public table and the evaluation function are identical.
The adversary view does not change.
No security argument depends on this change.

\subsection{Speedup}
The Lean median fell from 989 ms to 160 ms.
This is a 6.18 times speedup.
It is an 83.8 percent reduction.
This measurement uses the original 92-row incomplete law, before Sections~\ref{sec:o2} and~\ref{sec:o3}.

\section{Four-corner 91-digit recoding}
\label{sec:o2}

\subsection{Specification change}
Yes.
The scalar decomposition uses 91 fixed positions instead of 92.
The digit alphabet, the endomorphism, and the input-label shape do not change.
The Rust constant \texttt{KAPPA} changes from 92 to 91.
Old artifacts are not wire-compatible.

\subsection{Detail}
Let $r$ be the BN254 scalar modulus.
Let $\omega$ be the nontrivial cube root used by the G1 endomorphism.
Set $\beta=2-\omega$.
The digit alphabet is
\[
  D=\{0,\pm1,\pm\omega,\pm\omega^2\}.
\]
The norm of $\beta$ is seven.
Use the Eisenstein norm
\[
  Q(a,b)=a^2-ab+b^2.
\]
The relevant BN254 constants are
\begin{align*}
 A &=147946756881789319000765030803803410728,\\
 d &=9931322734385697763,\\
 C &=A+d,\\
 r &=A^2+Ad+d^2.
\end{align*}
The two GLV kernel vectors are
\[
  v_1=(-A,d),\qquad v_2=(-d,-C).
\]
Arkworks divides the two negative quotient numerators toward zero.
Thus, its GLV state is the floor representative $z_0$ used by the proof.
The selector checks
\[
  z_0,\quad z_0+v_1,\quad z_0+v_2,\quad z_0+v_1+v_2.
\]
Each shift is in the reconstruction kernel, so every candidate represents the same scalar.
For quotient remainders $u,v\in[0,r)$, scaling maps these four states to the four corners of one Eisenstein cell.
Two triangle lemmas show that one corner satisfies $3Q(a,b)\le r$.
This is the exact covering-radius step.

One recurrence round selects a unit digit and divides by $\beta$.
For a state with $Q(a,b)\le m$, the selected digit gives
\[
  Q(a-u_d,b-v_d)\le m+1+\lfloor\sqrt{4m}\rfloor,
\]
and exact division gives $Q(a-u_d,b-v_d)=7Q(a',b')$.
Define the inverse safe bound
\begin{align*}
 B(0)&=2,\\
 B(n+1)&=7B(n)+1-\left\lfloor\sqrt{4(7B(n)+1)}\right\rfloor.
\end{align*}
Lean proves $Q(a,b)\le B(n)\Longrightarrow \mathsf{after}(n+1,(a,b))=(0,0)$.
It also checks $r\le3B(90)$ with ten-round kernel-checked checkpoints and no \texttt{native\_decide}.
The main Lean results are:
\begin{itemize}
  \item \texttt{shiftedInitialCorrect};
  \item \texttt{existsCandidateNormBound};
  \item \texttt{fourCandidatesSignedI128};
  \item \texttt{existsCorrectShiftedInitialTerminates91};
  \item \texttt{shortInitialTerminates91};
  \item \texttt{shortInitialCorrect}.
\end{itemize}

Lean also checks $7^{90}<r<7^{91}$.
A fixed 90-position string over seven symbols has at most $7^{90}$ values, so 91 is optimal for this alphabet.

\subsection{Security analysis}
Every scalar uses exactly 91 positions.
Thus, the public circuit has 91 output rows for every scalar, and the topology does not reveal the scalar.
The Rust selector always runs all four 91-round tests.
It selects the first terminating candidate with \texttt{subtle::Choice} and has no secret-dependent early return.
Lean proves that every tested coordinate fits a signed 128-bit integer, so the fixed-width arithmetic cannot overflow.
The report does not claim compiler-level constant time.
The recoding is a public deterministic function of the secret scalar with the same alphabet, so the hybrid proof structure of the paper does not change.

\subsection{Speedup}
The dominant layer has 91 rows instead of 92.
This removes one row of 13 digit adaptors and 3,302 bit adaptors.
The per-row cost is unchanged.
The change also removes a sampled termination claim from the trust base and replaces it with a kernel-checked proof.

\section{Complete homogeneous addition}
\label{sec:o3}

\subsection{Specification change}
Yes.
The point row uses Algorithm~7 of Renes, Costello, and Batina in homogeneous coordinates instead of the exceptional affine law.
The row polynomials, the coefficient support flags, and the decode rule change.
The Rust constant \texttt{PRF\_VERSION} rises to 6 and binds the law into the setup commitment.

\subsection{Detail}
Let the hidden offset be $K=(A:B:C)$.
Let $(u,v)$ be the selected endomorphism image of the evaluator input.
Use homogeneous coordinates for $Y^2Z=X^3+3Z^3$.
An affine point is $(x:y:1)$ and the identity is $(0:1:0)$.
The three row polynomials are
\begin{align*}
 X={}&-9AB-18BCu-18ACv+(B^2-9C^2)uv+ABv^2,\\
 Y={}&-81C^2+27A^2u+27ACu^2+B^2v^2,\\
 Z={}&9BC+(B^2+9C^2)v+BCv^2+3A^2uv+3ABu^2.
\end{align*}
One fresh nonzero value $\rho$ multiplies all three coordinates.
The zero digit returns $\rho(A,B,C)$.
The decoder keeps $Z$ until the final projection.
For $Z\ne0$ it validates $(X/Z,Y/Z)$.
For $Z=0$ it accepts only $X=0$ and $Y\ne0$ and returns the identity.
It rejects all other zero-denominator triples.

The old affine law fails perfect correctness.
For scalar 2, output row 1 has digit \texttt{.one}.
Select the valid input equal to that row offset.
The old denominator is zero, the unchecked affine conversion becomes $(0,0)$, and decoding fails.

Lean proves the homogeneous curve equation, that the output is nonzero and projectively nonsingular, and that projective decode equals Mathlib group addition for every valid affine pair.
The proof has separate cases for different x-coordinates, doubling, and inverse inputs.
The main results are \texttt{outputNonsingular}, \texttt{outputRepresentsSum}, \texttt{decodeHomogeneous\_eq\_toAffine}, \texttt{decodeEvaluateRowSome}, \texttt{decodeEvaluateRowNone}, and \texttt{perfectCorrectness}.

The complete Bosma--Lenstra family has projective parameters $[\alpha:\beta:\gamma]$ with $\beta\ne0$.
Normalize $\beta=1$.
Under the exact \texttt{Biquadratic} support rule the adaptor counts are $D(0,0)=13$, $D(0,\gamma)=14$ for $\gamma\ne0$, and $D(\alpha,\gamma)=15$ for $\alpha\ne0$.
Thus, the RCB parameters $[0:1:0]$ give the unique minimum of 13.
A second bound covers private role-preserving linear rows.
A row with $t$ scalar adaptors exposes $14+t$ selected field values and has adaptor randomness of dimension at most $2t-2$.
Only the three output coordinates can stay unmasked, so $2t-2\ge 11+t$ and $t\ge13$.
A correct ten-adaptor sharing attempt violates this bound and exposes four extra coefficient invariants.

\subsection{Security analysis}
The change repairs perfect correctness.
BN254 G1 has cofactor one and odd prime order, so it has no point of order two.
The RCB exceptional line $Y=0$ does not meet this group.
Thus, the law covers doubling, inverse addition, and identity inputs for all valid G1 points.
The common factor $\rho$ randomizes the homogeneous representation, so the row does not reveal the offset scale.
Every support flag is public and independent of the scalar.
The $Z=0$ branch of the decoder reveals only that the decoded output is the group identity.
The 13-adaptor row is optimal in the fixed bidegree $(2,2)$ family and in the private role-preserving linear model.
The row exposes no coefficient function beyond the three outputs.

\subsection{Speedup}
This change increases work.
Each row uses 13 digit adaptors instead of nine.
Table~\ref{tab:cost} gives the exact counts.
The Lean median rises from 160 ms to 231 ms.
The complete circuit remains 4.28 times faster than the 989 ms baseline.
The HMAC table grows by 2,970,240 bytes to 9,668,659 bytes.

\begin{table}[h]
\centering
\begin{tabular}{lrrr}
\toprule
Item & Original 92-row call & Complete 91-row call & Change \\
\midrule
Output MAC rows & 92 & 91 & $-1$ \\
Digit adaptors & 833 & 1,188 & $+355$ \\
Bit adaptors & 211,582 & 301,752 & $+90,170$ \\
Fixed-permutation applications & 6,322,060 & 1,508,760 & $-4,813,300$ \\
Public permutation instances & 12,495 & 17,820 & $+5,325$ \\
\bottomrule
\end{tabular}
\caption{Circuit cost. The first column counts repeated calls in the unoptimized baseline. The last column uses one shared pass for each bit and point row.}
\label{tab:cost}
\end{table}

\section{Parallel rows with fixed RNG streams}
\label{sec:o4}

\subsection{Specification change}
Yes, in the randomness derivation.
The wire format does not change.
The garbler samples one fresh secret 32-byte root from the caller RNG.
Row $j$ uses ChaCha12 stream
\[
  \mathtt{0x666d3265635f6d61}+j
\]
from word position zero.
The code allocates the three child nonce domains of every row before Rayon starts.
The specification records this schedule.

\subsection{Detail}
The Rust implementation allocates one job for each of the 91 output rows.
Release builds execute the jobs with Rayon.
Debug builds use the same RNG and nonce schedule in serial order so the label-reuse diagnostic stays complete.
Indexed collection keeps the row order.
Jobs borrow the output keys and do not copy secret keys into a second vector.
The root seed is erased after row construction.
Fixed Boolean arrays replace heap-based coefficient index sets; this removes 27 net lines and 273 tree constructions per call.

\subsection{Security analysis}
Worker order cannot change the nonce domains, the caller RNG state, the row randomness, or the serialized row order.
Tests compare one-worker and four-worker tables and the final caller RNG state.
A second test compares all 273 allocated child nonces with the serial schedule.
The field-row seed is secret.
The pair of field-row seed and stream number must not repeat, and the security policy now states this requirement next to the existing seed, pre-seed, and nonce requirements.
Under this requirement the row randomness is a fresh pseudorandom stream per row, which is the same assumption the serial code already makes on its single stream.

\subsection{Speedup}
The serial 92-row incomplete baseline row estimate is 15.934 ms.
The complete 91-row, ten-worker row mean after this change is 3.845 ms.
This is 4.14 times faster.
A shared base row RNG did not improve this time and was discarded.

\section{Fixed AES window view reuse}
\label{sec:o5}

\subsection{Specification change}
None.
The fixed AES keys, inputs, and outputs do not change.

\subsection{Detail}
Each digit adaptor creates one fixed AES window view for each 91-bit chunk instead of one view per bit.
The view holds the 15 expanded keys of that window.

\subsection{Security analysis}
The change reorders allocation only.
The same keys encrypt the same labels.

\subsection{Speedup}
The measured row mean changed from 3.845 ms to 3.841 ms.
This is within noise.
A full AES block batching variant was slower at 4.578 ms and was discarded.

\section{SHA-256 prefix reuse}
\label{sec:o6}

\subsection{Specification change}
None.
Every derived key byte is identical.

\subsection{Detail}
The key derivation for the 15 fixed AES keys of one window shares one SHA-256 prefix for the bucket and role fields.
The code hashes the prefix once and forks the state for each key.

\subsection{Security analysis}
A test compares every derived cipher with the original full KDF input.
The output is the same function of the same input, so the key distribution does not change.

\subsection{Speedup}
The row mean fell from 3.841 ms to 3.754 ms.
This is a 2.3 percent reduction.

\section{Specialized hash-to-field reduction}
\label{sec:o7}

\subsection{Specification change}
None.
The 384-bit integer and its reduction modulo $p$ do not change.

\subsection{Detail}
The reducer assembles the fixed 48-byte hash input as a 17-byte low part and a 31-byte high part directly into field limbs.
It replaces the generic slice loop.

\subsection{Security analysis}
All 384 one-bit inputs and the prior edge vectors match the Arkworks reduction.
A compile-time constant for $2^{136}$ in $\mathbb{F}_q$ did not help and was discarded.

\subsection{Speedup}
The one-worker row mean fell from 20.858 ms to 19.906 ms.
This is a 4.6 percent reduction.
The ten-worker row mean is 3.630 ms.

\section{Four-limb field XOR}
\label{sec:o8}

\subsection{Specification change}
None.

\subsection{Detail}
The 256-bit field plaintext is XORed into its pad as four explicit little-endian 64-bit limbs.
This replaces nested base-field, limb, and byte loops.

\subsection{Security analysis}
XOR is bitwise, so the limb order changes no output bit.
The ciphertext row is the same value.

\subsection{Speedup}
The row mean fell from 3.713 ms to 3.390 ms.
Criterion reports an 8.7 percent improvement.
A native-endian 128-bit Davies--Meyer feed-forward XOR was 5.5 percent slower and was discarded.

\section{Direct hash-block reduction}
\label{sec:o9}

\subsection{Specification change}
None.

\subsection{Detail}
The three AES hash blocks feed the field reduction directly.
The code no longer copies them into a 48-byte staging array first.

\subsection{Security analysis}
The specialized limb split preserves every bit of the previous little-endian integer.
The hash-to-field value is identical.

\subsection{Speedup}
The row mean fell from 3.507 ms to 3.351 ms.
Criterion reports a 4.5 percent improvement.
The one-worker row mean is 18.411 ms.

\section{One-block KDF input}
\label{sec:o10}

\subsection{Specification change}
Yes.
The fixed-key KDF input packs the level, sub-bucket, role, and slot into one injective byte.
The complete input is 54 bytes.
\texttt{PRF\_VERSION} 5 binds the new domain and schedule into the setup commitment.

\subsection{Detail}
The 54-byte input and its SHA-256 padding fit in one compression block.
The metadata byte packs the level into bits 0--2 and the window-major key index into bits 3--5.
A contiguous 54-byte buffer did not improve the benchmark further and was discarded.

\subsection{Security analysis}
The valid ranges make each digest input distinct.
Domain separation is preserved because the version tag and the injective metadata map distinct keys to distinct random-oracle inputs.
The KDF is modeled as a random oracle, so each key is an independent uniform value as before.

\subsection{Speedup}
The row mean fell from 3.375 ms to 3.081 ms.
This is an 8.7 percent reduction.
The one-worker row mean is 16.950 ms.

\section{Paired KDF digests}
\label{sec:o11}

\subsection{Specification change}
Yes.
One SHA-256 digest supplies two 128-bit AES keys.
Eight digests derive the 15 keys of one window.
\texttt{PRF\_VERSION} 6 and the domain tag \texttt{afk/6} bind this schedule.

\subsection{Detail}
For flat key index $i$, set $j=\lfloor i/2\rfloor$ and $h=i\bmod2$.
Key $i$ is the $h$-th 128-bit half of digest $j$.
The second half of the last digest is unused.
Skipping its AES key expansion did not change the benchmark and was not adopted.

\subsection{Security analysis}
The map from each flat key index to its digest pair and half is injective.
In the random-oracle model both 128-bit halves of each 256-bit digest are independent uniform values.
Thus, the 15 keys remain independent and uniform.

\subsection{Speedup}
The row mean fell from 3.121 ms to 2.876 ms.
Criterion reports a 7.9 percent improvement.
The one-worker row mean is 15.348 ms.

\section{Noncanonical coordinate rejection}
\label{sec:o12}

\subsection{Specification change}
Yes, in the evaluator acceptance rule.
The checked evaluator rejects a 254-bit coordinate half that is not less than $p$ before field reduction.
It then checks the curve equation before it evaluates the HMAC table.

\subsection{Detail}
A regression test uses $x=p+1$ and $y=2$.
Reduction would otherwise map this input to the valid generator $(1,2)$.
The Lean public wrapper uses the canonical \texttt{AffineInput} and proves that its 508 encoding labels match \texttt{affineLamportBits} in x-then-y order.

\subsection{Security analysis}
The signed raw bits, the proof-derived canonical binding, the checked gadget path, and the final message-hash check must all agree for decryption to succeed.
A noncanonical encoding can no longer alias a different valid point.

\subsection{Speedup}
None.
The change is a hardening.

\section{End-to-end results}

The complete row stage is 5.54 times faster than the serial baseline: 2.876 ms against 15.934 ms.
The 95 percent interval of the final row mean is 2.861 to 2.893 ms.
The protocol \texttt{encode} mean is 68.190 ms with a 95 percent interval of 67.963 to 68.581 ms.
The comparable baseline is 81.309 ms.
The complete implementation reduces this time by 16.1 percent.
The Lean direct-table median is 231 ms against 989 ms, which is 4.28 times faster.

\section{Formal verification status}
\label{sec:formal}

The following invariants stay fixed:
\begin{itemize}
  \item Every candidate reconstructs the same scalar modulo $r$.
  \item Every public circuit has exactly 91 output positions.
  \item The digit alphabet does not change.
  \item The input keeps 508 Lamport label positions.
  \item The fixed-key window loads stay 91, 91, and 72.
  \item Every support flag is public and independent of the scalar.
  \item One nonzero factor scales all three homogeneous coordinates.
  \item Each row uses a separate nonce subtree and RNG stream.
  \item Worker order cannot change the table or caller RNG state.
\end{itemize}

The Rust selector uses four fixed 91-round tests and \texttt{subtle::Choice}.
This gives a fixed iteration count.
Lean proves that every tested coordinate fits a signed 128-bit integer.
The report does not claim compiler-level constant time.

The Rust checked evaluator rejects a coordinate that is at least the base-field modulus before reduction.
It then checks the curve equation.
A regression test uses $x=p+1$ and $y=2$.
Reduction would otherwise map this input to the valid generator $(1,2)$.

Lean proves that the encoding key contains the 508 Lamport label pairs in x-then-y order.
It proves that encoding selects the label for each bit of \texttt{affineLamportBits}.
The theorem is \texttt{Lamport.compatible}.

Lean proves exact joint-tape laws for fresh permutation and random-oracle programming.
Each map sends $(O,t)$ to the programmed oracle and the old value $O(x)$.
Each map is an involution on a finite sample space.
Thus, each map preserves the uniform joint PMF exactly.
Programming a fixed output alone does not preserve uniformity.
The permutation law also composes over every fixed finite target-tape schedule.
Both oracle laws preserve the uniform fiber of fresh oracles that match a fixed query transcript.
Lean lifts an exact handler equivalence through every bounded adversary oracle program.
Exact replay transports each real trace to the recording handler and preserves its result, query list, count, and final state.
A shared-sample coupling gives both traced marginals, equal results and query lists, and related final states.
Its structural form composes both adversary phases through an exact related-state bridge with zero result-disagreement mass.
The bounded form accepts a bridge that returns an exact result or records a bad event.
It keeps both marginals exact by using an independent coupling only after a bad event.
Lean proves that every result disagreement implies the bad flag.
It also proves that the final bad mass equals the bridge bad probability.
The final coupling theorem bounds the result advantage by that probability.
The permutation and hash involutions extend to the complete shared simulator state.
Fresh state swaps preserve every shared transcript and label invariant.
Each safe ideal-oracle query commutes with the matching full-state swap.
Thus, every path-safe bounded adversary program commutes with either swap in exact PMF equality.
Lean proves the exact mass of every finite forbidden block set.
A set with at most $k$ values has mass at most $k\cdot 2^{-128}$.
The two-point collision event therefore has mass at most $2^{-127}$.
Lean records each reached adversary query and its state.
The trace length is at most its indexed query budget.
Traced real and ideal games erase exactly to the challenge PMFs.
Each gives zero mass to paths above the sum of both adversary budgets.
If each step forbids at most $k$ blocks, its collision schedule has mass at most $qk\cdot 2^{-128}$.
An exact coupling lemma converts disagreement-event mass to the challenge's real-valued advantage.
A graph coupling maps each real trace to one ideal trace and proves both marginals.
A final generic theorem reduces the complete adaptive property to this trace transport and its Boolean change event.
Lean also bounds this change event by any containing reached-collision event.
The gate programmer accepts an ordered selected-gate schedule.
It preserves the shared invariant and any prior collision flag.
Every reached step is fresh unless the final state records a collision.
It counts all selected slots exactly.
A false-bit gate adds three fixed-key records, and a true-bit gate adds two.
A collision-free schedule adds their sum, which is at most three times the gate count.
The final oracle also satisfies every selected target after all later gate programs.
Lean instantiates this schedule for one complete 254-bit digit adaptor.
It proves the requested weighted field result after collision-free programming.
Lean also instantiates one schedule for all 13 active adaptors of one complete RCB row.
Its 3,302 selected bit gates return all three requested coordinate values in the same final oracle.
It then composes the five membership adaptors with all 91 RCB rows.
The complete fixed-key schedule has 1,188 digit adaptors and 301,752 selected bit gates.
The linked schedule derives its point-layer input MAC with the exact EncPRF transform.
It uses the programmed membership result as the whitening-key input.
Fixed-key programming preserves the EncPRF and hash oracles.
A satisfied linked schedule equals the full pipeline evaluation and returns all requested values.
The programming bridge now retargets each selected curve and RCB coordinate.
It puts one field residual in the low bit position and puts zero in every other position.
Lean proves that this weighted fold equals the residual.
It also proves that retargeting preserves every public table.
Collision-free programming returns each requested homogeneous value for all 91 rows.
The ideal output vector puts the requested point in row zero and the identity in the other 90 rows.
Lean proves that each homogeneous row decodes correctly and that radix Horner evaluation returns the requested point.
Retargeting to this vector preserves the complete public point table.
Lean also derives a finite instance for the complete ArgoMAC randomness type.
This type includes every algebraic value, label, permutation oracle, and random oracle.
The uniform security tape has full support over every complete randomness value.
Lean now defines the concrete online simulator under the challenge interface.
It returns the selected labels, retargets the curve and point requests, and applies the exact linked schedule.
The offline simulator coin is finite and contains no scalar or selected input.
It samples every public coefficient, ciphertext row, label pair, permutation family, and random oracle uniformly.
An explicit zero-valued coin supplies the nonempty witness, so the simulator needs no witness from its caller.
Lean proves that a finite equivalence transports a uniform PMF between two finite types.
It applies this result to the three real curve coefficients.
For every fixed hidden context, these coefficients are an affine bijection of the bridge key and two uniform masks.
Lean proves affine bijections for all public RCB X, Y, and Z coefficient groups.
It also proves that the real \texttt{garbleX}, \texttt{garbleY}, and \texttt{garbleZ} tables equal these transports.
The common zero pad is the sum of the public \texttt{y10} and \texttt{x9} digit offsets.
The digit-adaptor offset is independent of its private slope.
The 384-bit hash domain splits into $p\lfloor 2^{384}/p\rfloor$ complete field fibers and one suffix.
Lean proves exact uniform field transport on the complete fibers.
The equivalence uses residue-major order.
Thus, its field coordinate equals the actual reduction modulo $p$.
The offline coin samples one independent quotient for each of the 301,752 selected gates.
Retargeting preserves each quotient and changes only its residue.
It also proves that the suffix has exact mass
$(2^{384}\bmod p)/2^{384}$, which is at most $p/2^{384}$.
Lean applies a finite union bound to all selected gates.
The aggregate suffix mass is at most $301752p/2^{384}$.
This proof needs no independence assumption.
Lean also converts this bound to the stated real-valued hash-lift rounding error.
Lean now proves the required local marginal for one real gate.
A three-step swap schedule extracts the outputs of the three distinct fixed-key permutations.
The schedule preserves the uniform joint oracle tape exactly.
Davies--Meyer feed-forward is a componentwise involution, and three blocks are equivalent to one 384-bit value.
Thus, the construction's 384-bit gate hash is exactly uniform.
Its rejected suffix has the exact local mass above.
Lean splits the complete garbling tape into a uniform fixed-key oracle and all remaining randomness.
A finite mixture theorem lifts the local law to the complete security tape.
The label can depend on every non-fixed-oracle value.
Lean now indexes the exact selected schedule by all 301,752 gates.
The metadata contains each real location, window, and selected curve or linked point label.
The full suffix union has mass at most $301752p/2^{384}$ on the complete tape.
Lean also proves the exact law for one actual bit-adaptor ciphertext row.
A two-step swap extracts both pad-permutation outputs and leaves every hash slot unchanged.
Thus, the real plaintext can contain its hash-to-field offset.
Davies--Meyer feed-forward, block joining, and XOR by the encoded field value are finite equivalences.
The true row is exactly uniform on 256 bits.
A finite mixture permits keys and slopes that depend on all non-fixed-oracle randomness.
Lean now proves the good-event law for the actual zero pad.
One finite product equivalence exposes each accepted hash lift as its exact residue and quotient.
The low digit has coefficient one, so its residue is an additive pivot.
Thus, a uniform accepted-residue vector gives an exact uniform digit offset after any fixed translation.
The actual good \texttt{y10} and \texttt{x9} hash vectors give the two real offsets.
Their sum is the real zero pad.
Lean now proves that one gate's three hash reads and two pad reads are jointly uniform.
The five fixed-key indices differ only in their slot field, so they are distinct.
A five-slot swap schedule extracts all five outputs and preserves the uniform tape.
Thus, the hash-based zero pad and the pad-based ciphertext row are independent on one gate.
Lean also proves the general distinct-read tool.
For any injective family of fixed-key indices, the joint permutation reads are uniform, with no bound on the count.
The proof shows by induction that a schedule with distinct slots and distinct indices reads each entry at its slot.
The window-sharing optimization reads one permutation at 91 labels per window, so the full read family is not index-injective.
The joint transport therefore factors as independence across the 17,820 distinct window indices, times the within-window collision that the existing bad event already bounds.
The next proof must package each window's read as one value, instantiate the distinct-read tool across windows, and compose the result with the coefficient transports and the collision bound.

The complete row repairs the known perfect-correctness failure.
The generic bounded bridge is complete.
The construction still lacks the joint distributional hop for the RCB zero pad and ciphertext rows.
It also lacks the final adaptive collision bound for that hop.
The worst selected branch uses the three hash slots.
Its zero-query numerator is
\[
  3\cdot 10692\cdot 91^2=265621356<2^{28}.
\]
The 100-bit margin is 2,814,100 numerator units.
A direct five-slot union bound exceeds $2^{28}$.
Lean proves the 100-bit arithmetic for each exclusive three-hash or two-pad branch.
For a bucket with $f$ false selections and $t$ true selections, Lean proves
$3f^2+2t^2\leq 3(f+t)^2$ and $3f+2t\leq 3(f+t)$.
Thus, mixed selected bits stay below the three-hash bucket budget.
The semantic reduction must still condition on the selected branch.
A companion Kriterion change closes the challenge interface defects.
It makes \texttt{Kriterion.Solution} a closed dependent bundle.
It quantifies every proof over the BN254 field and group certificates, which remain declared challenge assumptions.
The certificate-free benchmark garbler has a required equality to the proved scheme garbler.
It also defines \texttt{Kriterion.Benchmark.garble} at a fixed scalar, security parameter, and seed.
The security game now samples the complete finite random tape uniformly.
The fixed 256-bit seed selects only the reproducible benchmark tape.
This change permits exact uniform permutation and random-oracle proofs.
The challenge must still reject this entry until the construction-specific adaptive proof is complete.

\section{Reproduction}

The Lean proof uses Lean 4.24 and the pinned Mathlib checkout.
Run:
\begingroup\small
\begin{verbatim}
lake build Proof Tests
lake env lean entry/tests/Correctness.lean
lake env lean AxiomCheck.lean
lake exe bench
\end{verbatim}
\endgroup

The Rust validation uses Rust 1.94.0.
Run:
\begingroup\small
\begin{verbatim}
cargo fmt --all -- --check
cargo clippy --all-targets --all-features -- -D warnings
cargo test --all-features
cargo test --release --all-features
RAYON_NUM_THREADS=10 cargo bench --all-features \
  --bench field_mac_to_ec_mac -- 'field_mac_to_ec_mac/g1/garble'
RAYON_NUM_THREADS=10 cargo bench --all-features \
  --bench protocol_bench -- 'encode'
\end{verbatim}
\endgroup

\section*{References}

J.~Bosma and H.~W.~Lenstra, ``Complete systems of two addition laws for elliptic curves,'' \emph{Journal of Number Theory}, 1995.

J.~Renes, C.~Costello, and L.~Batina, ``Complete addition formulas for prime order elliptic curves,'' EUROCRYPT 2016.

RustCrypto, \emph{AES 0.8.4 documentation}, ARMv8 run-time detection section.

Arkworks, \emph{ark-ec 0.5.0}, GLV scalar decomposition implementation.

\end{document}
